На экране анализа доступен каталог с поиском, фильтрами и запуском обработки
отдельной записи или выбранной группы (рисунки~\ref{fig:proto1} и~\ref{fig:proto2}).
\begin{figure}[H]\centering
 \reportScreenshot{figures/01-analyze-audio.png}
 \caption{Каталог записей для запуска анализа}\label{fig:proto1}
\end{figure}
\begin{figure}[H]\centering
 \reportScreenshot{figures/02-analyze-audio.png}
 \caption{Выбор аудиозаписей для анализа}\label{fig:proto2}
\end{figure}

Состояние выполняемых заданий и сведения о конкретном запуске показаны
на рисунках~\ref{fig:proto3} и~\ref{fig:proto4}.
\begin{figure}[H]\centering
 \reportScreenshot{figures/03-analyze-audio-active-task.png}
 \caption{Список активных заданий анализа}\label{fig:proto3}
\end{figure}
\begin{figure}[H]\centering
 \reportScreenshot{figures/04-analyze-audio-active-task-detailed.jpg}
 \caption{Подробные сведения об активном задании}\label{fig:proto4}
\end{figure}

Для повторной обработки всего каталога предусмотрен отдельный сценарий
(рисунок~\ref{fig:proto5}). История запусков содержит статусы и сведения
о качестве результатов, а также предоставляет экспорт (рисунки~\ref{fig:proto6}
и~\ref{fig:proto7}).
\begin{figure}[H]\centering
 \reportScreenshot{figures/05-analyze-catalog.png}
 \caption{Запуск анализа каталога}\label{fig:proto5}
\end{figure}
\begin{figure}[H]\centering
 \reportScreenshot{figures/06-view-analisis-history.png}
 \caption{Просмотр истории анализов}\label{fig:proto6}
\end{figure}
\begin{figure}[H]\centering
 \reportScreenshot{figures/07-export-analysis-history.png}
 \caption{Экспорт истории анализов}\label{fig:proto7}
\end{figure}

Вход администратора показан на рисунке~\ref{fig:proto8}. Операция удаления
результата представлена последовательностью выбора, подтверждения и уведомления
(рисунки~\ref{fig:proto9}-\ref{fig:proto11}).
\begin{figure}[H]\centering
 \reportScreenshot{figures/08-admin-login.png}
 \caption{Вход администратора}\label{fig:proto8}
\end{figure}
\begin{figure}[H]\centering
 \reportScreenshot{figures/09-delete-analisis-result.png}
 \caption{Выбор результата для удаления}\label{fig:proto9}
\end{figure}
\begin{figure}[H]\centering
 \reportScreenshot{figures/10-delete-analisis-result.png}
 \caption{Подтверждение удаления результата}\label{fig:proto10}
\end{figure}
\begin{figure}[H]\centering
 \reportScreenshot{figures/11-delete-analisis-result-notifications.png}
 \caption{Уведомление об удалении результата}\label{fig:proto11}
\end{figure}

Результаты анализа и характеристики записи представлены на
рисунках~\ref{fig:proto12} и~\ref{fig:proto13}. Экран параметров модели
позволяет задать версию модели и пороги уверенности (рисунок~\ref{fig:proto14});
отмена задания показана на рисунке~\ref{fig:proto15}.
\begin{figure}[H]\centering
 \reportScreenshot{figures/12-view-result-of-analysis.png}
 \caption{Просмотр результата анализа}\label{fig:proto12}
\end{figure}
\begin{figure}[H]\centering
 \reportScreenshot{figures/13-view-result-of-analysis.png}
 \caption{Подробные характеристики результата анализа}\label{fig:proto13}
\end{figure}
\begin{figure}[H]\centering
 \reportScreenshot{figures/14-config-model-parameters.png}
 \caption{Настройка параметров модели и порогов}\label{fig:proto14}
\end{figure}
\begin{figure}[H]\centering
 \reportScreenshot{figures/15-cancel-analyze-task.png}
 \caption{Отмена задания анализа}\label{fig:proto15}
\end{figure}
